% =====================================================================
%  Sezione 2.11 — Definizione del compito ed etichettatura delle finestre
%  Da inserire in Cap2.tex al posto della vecchia sezione
%  "Finestre temporali e labeling" (parte relativa all'etichettatura).
% =====================================================================

\section{Definizione del compito ed etichettatura delle finestre}

Fino a questo punto il problema è stato descritto come «distinguere una caduta
da una normale attività quotidiana». È la formulazione corrente in letteratura,
ma nasconde un'ambiguità che, come mostrerò nel capitolo sperimentale, produce
l'effetto più grande misurato in tutto questo lavoro.

\subsection{Una caduta è composta da tre fasi, non da una}

Osservando le riprese di caduta del dataset si distinguono tre stati
fisicamente diversi, che l'annotazione per fotogramma di UP-Fall separa
esplicitamente:

\begin{itemize}
    \item \textbf{attività normale}: la persona è in piedi, cammina o si siede;
    \item \textbf{transizione}: dalla perdita di equilibrio all'arrivo a terra;
    \item \textbf{stato post-caduta}: la persona è a terra.
\end{itemize}

Le tre fasi hanno durate molto diverse. Misurando le cinquanta riprese di
caduta della Camera~1 si ottengono i valori riportati in
Tabella~\ref{tab:durate_fasi}: la transizione dura circa tre secondi, mentre lo
stato successivo dura il doppio, e nel funzionamento reale può durare minuti od
ore.

\begin{table}[ht]
\centering
\caption{Durata delle tre fasi, misurata sulle riprese di caduta di UP-Fall
(Camera~1, dieci soggetti).}
\label{tab:durate_fasi}
\begin{tabular}{lccc}
\hline
\textbf{Fase} & \textbf{Mediana} & \textbf{Decile 10\%} & \textbf{Decile 90\%} \\
\hline
In piedi, prima della caduta & 0,92 s & --- & --- \\
Transizione (caduta in corso) & 3,11 s & 2,05 s & 3,94 s \\
Stato post-caduta & 5,97 s & 4,82 s & 6,89 s \\
\hline
\end{tabular}
\end{table}

\subsection{Tre compiti diversi, non tre gradazioni dello stesso}

Da questa scomposizione discende che «rilevare una caduta» non individua un
compito solo. A seconda di quale fase si assume come positivo se ne definiscono
tre, con prevalenze, difficoltà e negativi critici diversi:

\begin{itemize}
    \item \textbf{evento}: positivo se la finestra contiene la transizione.
          Il sistema riconosce \emph{l'atto} del cadere;
    \item \textbf{stato}: positivo se la finestra contiene la persona a terra.
          Il sistema riconosce \emph{la condizione} di essere caduti;
    \item \textbf{allarme}: positivo se la finestra contiene l'una o l'altra.
\end{itemize}

La scelta fra i tre non è un dettaglio di preprocessing. Il criterio clinico
suggerirebbe il terzo, o addirittura il secondo: il danno della \emph{long-lie
syndrome} dipende dal tempo trascorso a terra, non dall'impatto, e lo stato dura
molto più a lungo della transizione, offrendo quindi più occasioni di
rilevamento. Questa è la formulazione che il presente lavoro ha adottato
inizialmente.

\subsection{Il negativo che rende la scelta non ovvia}

Esiste però una classe di negativi che rende i tre compiti profondamente
diversi in difficoltà: le persone che si trovano in posizione orizzontale
\emph{senza} essere cadute.

In UP-Fall una delle attività registrate come normali è \emph{laying}, cioè
sdraiarsi volontariamente: sono 20.980 fotogrammi, il 18\% dell'intero dataset,
etichettati come attività quotidiana. Il dataset contiene quindi ventimila
fotogrammi di persona distesa etichettati come normali e diecimila etichettati
come post-caduta. \textbf{Nello stato finale i due casi sono quasi
indistinguibili: la differenza sta unicamente nel modo in cui la persona ci è
arrivata.}

Questa osservazione è verificabile quantitativamente. Confrontando le mediane
delle grandezze dinamiche fra le popolazioni di fotogrammi (Tabella
\ref{tab:dinamica}), la transizione di caduta si separa dallo sdraiarsi
volontario di uno o due ordini di grandezza, mentre lo stato post-caduta e lo
stato disteso volontario sono praticamente sovrapposti.

\begin{table}[ht]
\centering
\caption{Rapporto fra le mediane delle grandezze dinamiche. La transizione
discrimina; la posa no.}
\label{tab:dinamica}
\begin{tabular}{lcc}
\hline
\textbf{Grandezza} & \textbf{Caduta / sdraiarsi} & \textbf{A terra / disteso} \\
\hline
Velocità del bacino & 30,1$\times$ & 1,3$\times$ \\
Escursione del tronco & 42,5$\times$ & 2,8$\times$ \\
Percorso del bacino & 33,0$\times$ & 1,3$\times$ \\
Discesa del bacino & 139,4$\times$ & 2,1$\times$ \\
\hline
\end{tabular}
\end{table}

La ragione è fisica e intuitiva: nessuno si sdraia su un divano con
un'escursione del tronco di trenta gradi in mezzo secondo. Lo stato finale è
ambiguo per costruzione, la transizione no.

\subsection{Dalle etichette per fotogramma a quella della finestra}

Il classificatore non decide sul singolo fotogramma ma su una finestra di circa
2,4 secondi, che quasi sempre contiene fotogrammi di fasi diverse. Serve quindi
una regola di aggregazione.

La regola adottata inizialmente contava i fotogrammi della classe ovunque
essi fossero nella finestra: cinque fotogrammi di caduta sparsi valevano quanto
cinque consecutivi. Per un evento che è per natura una dinamica continua il
criterio è incoerente. È stato quindi sostituito da un criterio
\textbf{$K$-su-$M$}: la finestra è positiva se esistono almeno $K$ fotogrammi
della classe all'interno di un blocco locale di $M$ fotogrammi consecutivi.

Formalmente, data la maschera binaria $m_t$ che vale 1 sui fotogrammi della
classe cercata, la finestra è positiva se

\[
\max_{i}\ \sum_{t=i}^{i+M-1} m_t \ \geq\ K .
\]

Il criterio tollera un fotogramma annotato male senza spezzare l'evento, e ha
due casi particolari utili: con $M = K$ degenera in «$K$ fotogrammi
consecutivi», con $M$ molto grande nel conteggio totale sulla finestra. Il
confronto fra criteri diventa quindi una variazione di parametro invece che di
codice.

I valori utilizzati sono $K = 5$ su $M = 8$ per la transizione e $K = 10$ su
$M = 15$ per lo stato. Non sono valori ottimizzati: sono soglie di
preprocessing scelte per imporre una presenza minima della fase nella finestra,
ed è dichiarata nel Capitolo~\ref{cap:sperimentale} un'analisi di sensibilità.

\subsection{Il metro di valutazione deve essere unico}

Confrontare tre definizioni del compito richiede una cautela: se ciascun
modello venisse misurato sulla propria definizione di positivo, il confronto non
sarebbe interpretabile. Il compito \emph{evento} scarta le finestre di stato,
che sono le più ambigue, e apparirebbe avvantaggiato per un motivo che non ha
nulla a che vedere con la sua qualità.

Ogni finestra porta quindi \textbf{due etichette}:

\begin{itemize}
    \item $y_{\text{train}}$, che dipende dalla definizione del compito ed è la
          variabile dell'esperimento;
    \item $y_{\text{allarme}}$, fissa per tutte le definizioni: vale 1 se la
          finestra contiene la transizione \emph{oppure} lo stato. È il metro
          con cui i tre modelli vengono confrontati.
\end{itemize}

Le finestre che il compito \emph{evento} esclude dall'addestramento restano
quindi nell'insieme di verifica, marcate come non addestrabili: tutti i modelli
sono misurati sullo stesso insieme.

\subsection{Un vocabolario di riferimento}

La tassonomia adottata coincide con quella di \emph{OmniFall}
\cite{omnifall2025}, un benchmark del 2026 che unifica otto dataset di
rilevamento cadute — UP-Fall compreso — per un totale di 101 soggetti e 29
punti di vista, riannotandoli a mano con sedici classi. Quattro di queste sono
esattamente le distinzioni discusse in questa sezione:

\begin{itemize}
    \item \texttt{fall} --- l'atto di cadere, da qualunque stato precedente;
    \item \texttt{fallen} --- essere a terra dopo una caduta;
    \item \texttt{lie\_down} --- sdraiarsi intenzionalmente;
    \item \texttt{lying} --- essere disteso dopo esserci andato intenzionalmente.
\end{itemize}

Gli autori distribuiscono le sole annotazioni e non i video, il che ha permesso
di applicarle ai dati già estratti senza modificare la pipeline: le etichette
temporali sono state mappate sui fotogrammi di UP-Fall con una copertura del
96,4\%, e l'allineamento è verificato dal fatto che le classi \texttt{fall} e
\texttt{fallen} cadono dove cadono le annotazioni originali del dataset.

L'adozione di questo vocabolario ha due conseguenze utili. La prima è che i
negativi critici diventano misurabili \emph{dentro} il dataset di
addestramento, e non soltanto sulle registrazioni domestiche o su un archivio
esterno. La seconda è che i risultati diventano confrontabili con un benchmark
pubblicato: OmniFall riporta, con un modello video da 87 milioni di parametri,
F1 pari a 65,6\% sul riconoscimento dell'evento e \textbf{41,0\% sul
riconoscimento dello stato}, e osserva che i modelli multimodali di grandi
dimensioni confondono sistematicamente \texttt{fallen} con \texttt{lying}. Il
divario fra i due compiti non è quindi un limite di questo lavoro: è un
problema aperto del settore.

\subsection{Cosa ne segue}

La definizione del compito è stata trattata in questo capitolo come una scelta
di rappresentazione del problema, alla pari della scelta delle feature, e non
come un dettaglio di preprocessing. Il Capitolo~\ref{cap:sperimentale} ne misura
le conseguenze, che risultano di gran lunga superiori a quelle della scelta
dell'architettura di classificazione.
